\section{Related Work}
\subsection{Learning chess evaluations from demonstrations}
Learning evaluation parameters from played moves predates the present study. David et al.\ \cite{david} evolved evaluation-function parameters to match moves from grandmaster games, then used coevolution for further improvement. Their approach is relevant because it uses demonstrated choices without requiring a numerical evaluation label for each training state. Thus, neither learning from moves nor reporting material parameters is, by itself, a novelty claim for our work. The current study instead uses a small probabilistic choice model and focuses on the interpretation and uncertainty of coefficients across explicitly defined position subsets.

DeepChess \cite{deepchess} learns to compare positions using neural representations trained from game data and integrates the learned comparison with search. Its treatment of evaluation illustrates a different path from fixed hand-specified features to learned representations. Our model deliberately retains a small explicit feature set because each material term must have a directly inspectable coefficient. This sacrifices expressive capacity and does not make our representation a replacement for a neural evaluator. The appropriate comparison is between the questions being answered: learning useful position comparisons versus auditing the material terms of a restricted behavioral reward model.

\subsection{Playing strength and behavior prediction}
AlphaZero \cite{alphazero} learns policy and value predictions through self-play and uses them to guide search. That forward learning problem asks how to generate strong behavior from an outcome-based objective. Our inverse problem conditions on behavior already generated by a fixed engine and estimates a compact reward description of the observed choices. We do not train an AlphaZero-style agent, infer its representations, or compare playing strength. A demonstrated move is the label here; the game result is retained as part of the dataset history rather than used as a material-value target.

Maia, introduced by McIlroy-Young et al.\ \cite{maia}, models human move decisions and shows why agreement with human choices is a separate objective from maximizing chess strength. This distinction is useful for interpreting our engine-based dataset as well. Matching an engine's decisions under a fixed search budget does not establish alignment with human judgments or explain why humans assign conventional values to pieces. Our source of demonstrations is therefore part of the estimand, not a neutral implementation detail. Extending the study to human games would require a separate population definition and evaluation rather than reinterpreting the existing coefficients as human preferences.

\subsection{Contextual material values and neural interpretation}
P\'alsson and Bj\"ornsson \cite{palsson} probe concepts learned by a strong chess-playing agent and study relative material values across game phases. Their work establishes a close substantive connection: contextual material evaluation is already an object of empirical investigation. Our analysis differs in the observation used to fit the model. Neural probing examines an existing agent's representations or evaluation behavior, whereas our coefficients are fitted from selected legal actions under a specified probabilistic reward model. Similar-looking value tables can therefore describe different properties of an agent.

PAWN \cite{pawn} predicts position-specific piece values using engine-derived labels associated with removing a piece. Such labels ask how an evaluator responds to a board modification. Our supervision consists only of the chosen move among legal alternatives, and our material features are differences induced by those alternatives. Removing a piece for an evaluation experiment and observing a legal capture need not produce equivalent statistical information. A removal-based target can also be available for a piece that is not immediately capturable, whereas a one-step material coefficient is constrained only when candidate actions differ in that feature.

The distinction is consequential for sparse boards. A state may contain a strategically essential minor piece but offer no immediate action that changes either side's count of that piece. Our immediate material feature then contributes the same value to every legal action and cannot affect their probabilities at that state. A neural evaluation probe can still respond to that piece's presence or location. Consequently, a discrepancy between our ratios and a probing or removal-based study need not be a failed replication. It may instead reflect a different target and a different set of informative observations. We make this issue measurable by reporting feature-specific support counts.

\subsection{Inverse reinforcement learning and maximum entropy}
Ng and Russell \cite{ng2000} formalized algorithms for recovering rewards from observed behavior, motivating the distinction between an observed policy and a reward that can explain it. Abbeel and Ng \cite{abbeel} developed apprenticeship learning through matching feature expectations and showed why useful behavioral performance need not require exact recovery of the expert's true reward. These ideas support an important separation in our study: a fitted reward can be a useful description without being a unique reconstruction of an underlying objective.

Ziebart et al.\ \cite{ziebart} introduced a maximum-entropy formulation that places a globally normalized distribution over decision sequences. Our legal-action softmax is a horizon-one specialization. The normalization is over the legal alternatives of one observed state, and fitting does not invoke a sequential chess planner. This is computationally convenient, but it changes the inference target. We explicitly use the phrase one-step maximum-entropy IRL to distinguish the implemented estimator from trajectory-level methods whose feature expectations depend on longer-term policies.

The contextual models in this paper alter the linear score's feature design. Global, phase-only, openness-only, and joint specifications remain members of the same maximum-entropy conditional-choice family. Independent group fits alter which observations share a parameter vector. These are meaningful estimation comparisons, but they do not constitute an algorithm benchmark. For example, generative adversarial imitation learning \cite{gail} learns behavior through an adversarial framework; no such method is executed here. Its inclusion in the conceptual landscape should not be confused with an experimental baseline.

\subsection{Reward identification and coefficient interpretation}
Kim et al.\ \cite{kim} study reward identification in IRL and characterize conditions under which rewards can be identified. Their work cautions against equating a fitted optimum with unique knowledge of the demonstrator's objective. In our restricted conditional model, ridge regularization makes the numerical optimization well behaved, but that property does not resolve the broader identification problem or compensate for omitted tactical continuations. We distinguish uniqueness within a fixed penalized parameterization from recovery of an engine's strategic reward.

\subsection{Positioning the present study}
The paper connects a deliberately transparent one-step estimator to two concrete partitions of a fixed chess dataset: operational game stages and remaining-piece counts. Its contribution is an application and audit of a restricted model, with uncertainty and representation dependence made explicit. It does not establish that contextual piece values are newly discovered, that IRL was previously unavailable for demonstrations, or that a small linear model captures all of chess strategy. The new experimental extension is the actual independent fitting across piece-count groups, with shared-game resampling and support diagnostics reported alongside the original stage analysis.